Dues partícules carregades es mouen en el pla del paper a la mateixa velocitat per una zona en què hi ha un camp magnètic uniforme de valor $4,50\times 10^{-1}\ \mathrm{T}$ perpendicular al pla i que surt del paper (vegeu la figura). Part de les trajectòries descrites per les càrregues són les que es veuen també en la figura. La partícula $Q_1$ té una massa de $5,32\times 10^{-26}\ \mathrm{kg}$ i la partícula $Q_2$, de $1,73\times 10^{-25}\ \mathrm{kg}$. La magnitud de cadascuna de les càrregues és la mateixa, $3,20\times 10^{-19}\ \mathrm{C}$, i la força magnètica que actua sobre elles també té el mateix mòdul, que és $1,01\times 10^{-12}\ \mathrm{N}$.

\figura[0.3]{fig1.png}

\begin{apartat}{1}
Expliqueu raonadament el signe que tindrà cadascuna de les càrregues. Calculeu la velocitat d'aquestes càrregues.
\end{apartat}

\begin{apartat}{1}
Calculeu els radis de les trajectòries de cada partícula i la freqüència (Hz) del moviment de $Q_2$.
\end{apartat}
